\section{Relation to earlier work}\label{sec:related}

Bukh and Cox~\cite[Theorems~2(b) and~3(b)]{BC} proved a bound for general
matrices. For a complex matrix $A$ of order
$N$ with diagonal one and rank at most $N-k$, it gives
\begin{equation}\label{eq:bc-literature}
 \max_{i\ne j}|A_{ij}|\geq\frac{1}{\alpha_kN-1},
 \qquad
 \alpha_k=\frac{1+(k-1)\sqrt{k+1}}{k^2}.
\end{equation}
Although their extremal quantity is defined using exact rank, the proof
of their Theorem~10 already permits rank at most $N-k$. Equality in
\eqref{eq:bc-literature} requires a maximal equiangular tight frame in
$\mathbb C^k$ and $k^2\mid N$; uniform replication supplies the matching
construction. In codimension two this is the tetrahedral configuration
of four complex lines. Consequently, the exact value
$((1+\sqrt3)N/4-1)^{-1}$ when $4\mid N$ is an established result.

Their corrigendum~\cite{BCErratum} specifies the complex equality case:
equal mass is assigned to each entire
unit-phase orbit, rather than to a pair of antipodal representatives.
This correction preserves the numerical bound and its divisibility
condition. We use the corrected projective formulation throughout.

Magsino, Mixon and Parshall~\cite[Lemma~1 and Theorem~4]{MMP} give a
Delsarte-style proof through an auxiliary polynomial and tight-frame
duality. Glazyrin~\cite[Theorem~7]{Glazyrin} develops a general matrix
duality for conjugate energies. His replicated-code construction and
equality analysis~\cite[Theorem~10 and the subsequent discussion]{Glazyrin}
also show that an extremizing general matrix is a Gram matrix, with
uniqueness inherited from the maximal projective simplex. In particular,
the divisible codimension-two case already has this stronger equality
description. We use this duality, the auxiliary-polynomial estimate,
phase alignment at equality, and the tetrahedral model in the proof below.

Two closely related optimization problems help explain the dual
functional. Cahill and Casazza~\cite[Proposition~2.1 and Theorems~2.2,
2.4]{CahillCasazza} study total coherence of Parseval frames, its
invariance under Naimark complementation, and restrictions on maximizing
frames. For columns $u_i$ of a matrix $U\in\mathbb C^{2\times N}$ with
$UU^*=I_2$, write
\[
 \operatorname{TC}(U)=\sum_{i\ne j}|u_i^*u_j|.
\]
The quasimaximal relative projection constant of Foucart and
Skrzypek~\cite[Section~2]{FoucartSkrzypek} is, in this notation,
\[
 q_{\mathbb C}(2,N)
 =\frac{\max_{UU^*=I_2}\operatorname{TC}(U)+2}{N}.
\]
Their maximal relative projection constant additionally optimizes
independent nonnegative weights $t_i$, with $\sum_i t_i^2=1$, in
$\sum_{i,j}t_it_j|u_i^*u_j|$. These distinctions matter when transferring
an exact value or an equality statement to the coherence problem.

Kobos~\cite[Theorem~4.2]{Kobos} proves a stability theorem: a sequence
approaching the maximal ETF bound admits, after
passing to a subsequence, clustering around maximal ETF directions in
a weighted mean-square sense; an exceptional set may have vanishing
total weight. His analysis also treats phases and residuals. To obtain
the finite conclusion here, we need bounds that place every column in
a controlled neighborhood at each cardinality under consideration.
The subsequent local comparison must eliminate all within-cluster
residuals. These are the quantitative steps developed in this paper.

Der\k{e}gowska and Lewandowska~\cite[Theorem~3.6]{DLWelch} determine
relative projection constants at replicated maximal-ETF cardinalities.
Kumar, Mohar, Mojallal and Pragada~\cite[Theorem~1.3]{KMMP} establish
eventual stabilization of maximal relative projection constants in the
real setting. The codimension-two theorem of Der\k{e}gowska and
Lewandowska~\cite[Theorem~2.4]{DLCodim} instead controls a real
Banach--Mazur distance through a small projection constant. These
results concern different objectives from the finite complex
off-diagonal minimax problem in Theorem~\ref{thm:main}.

Theorem~\ref{thm:main} extends the exact value and complete equality
classification to the nondivisible cardinalities in the range
$N\geq180$, within the full class of complex unit-diagonal matrices of
rank at most $N-2$. It identifies the balanced integer multiplicities,
the corresponding deformation of the four-line configuration, and all
extremizers. Quantitative local rigidity and computer-assisted existence
arguments also occur in projective coding theory, for example
in~\cite{CKM}. Here the local comparison is combined with control of
every kernel column to obtain a global minimax theorem uniformly over
the stated cardinalities.
